% !TEX root = main.tex
% ===============================================================
%  boyogalazhcnbooksetting.tex —— boyogalazhcnbook 中文书籍模板 全局设置文件
%  由 main.tex 自动载入 (\input); 标题/作者等封面信息仍在 main.tex 中设置.
%  所有字体均取自本仓库 font/ 目录, 在任何电脑上编译都无需安装字体.
% ===============================================================
% ===============================================================
%  一、字体设置
%  模板自带完整字体方案: 全部字体取自仓库 font/ 目录, 不依赖系统字体,
%  也不使用任何第三方字体集; 中文命令与字体族均在本节定义 —
%    命令: \songti(宋) \heiti(黑) \kaishu(楷) \fangsong(仿) \coverfont(封面)
%    字族: cnsong / cnhei / cnkai / cnfs / cncover
% ===============================================================

%----- 中文字体: Fandol 字族 (取自 font/) -----
% 文类以 fontset=none 载入, 中文完全在此设置; 字体族取名 cnsong/cnhei/cnfs/cnkai,
% 不与 ctex 预置名冲突 (避免字体重定义警告), 字体文件均取自 font/ 目录.
\setCJKmainfont{FandolSong-Regular.otf}[        % 正文: 宋体
    Path = font/,
    BoldFont = FandolSong-Bold.otf,             % 粗体: 宋体粗
    ItalicFont = FandolKai-Regular.otf]         % 斜体: 楷体
\setCJKsansfont{FandolHei-Regular.otf}[         % 无衬线: 黑体
    Path = font/,
    BoldFont = FandolHei-Bold.otf]
\setCJKmonofont{FandolFang-Regular.otf}[        % 等宽: 仿宋
    Path = font/]
% 中文字体族与中文命令: \songti \heiti \fangsong \kaishu 均指向本地字体文件
% (fontset=none 时 ctex 不提供这四个命令, 用 \providecommand + \renewcommand 定义)
\setCJKfamilyfont{cnsong}{FandolSong-Regular.otf}[Path = font/, BoldFont = FandolSong-Bold.otf]
\setCJKfamilyfont{cnhei}{FandolHei-Regular.otf}[Path = font/, BoldFont = FandolHei-Bold.otf]
\setCJKfamilyfont{cnfs}{FandolFang-Regular.otf}[Path = font/]
\setCJKfamilyfont{cnkai}{FandolKai-Regular.otf}[Path = font/]
\providecommand\songti{}\renewcommand\songti{\CJKfamily{cnsong}}
\providecommand\heiti{}\renewcommand\heiti{\CJKfamily{cnhei}}
\providecommand\fangsong{}\renewcommand\fangsong{\CJKfamily{cnfs}}
\providecommand\kaishu{}\renewcommand\kaishu{\CJKfamily{cnkai}}

%----- 封面字体: 思源黑体 (展示用, 比 Fandol 黑体更现代; 供文类封面使用) -----
\setCJKfamilyfont{cncover}{SourceHanSansCN-Regular.otf}[Path = font/, BoldFont = SourceHanSansCN-Bold.otf]
\renewcommand{\coverfont}{\CJKfamily{cncover}}   % 封面标题/副标题/作者均使用该字体

%----- 西文字体 -----
% 衬线 (正文): STIX Two Text —— Times 字体现代修订版, 学术出版风格, 与数学字体同族
\setmainfont{STIXTwoText-Regular.ttf}[
    Path = font/,
    BoldFont = STIXTwoText-Bold.ttf,
    ItalicFont = STIXTwoText-Italic.ttf,
    BoldItalicFont = STIXTwoText-BoldItalic.ttf]
% 无衬线: 思源黑体 (其西文字形为 Source Sans, 干净现代)
\setsansfont{SourceHanSansCN-Regular.otf}[
    Path = font/, BoldFont = SourceHanSansCN-Bold.otf]
% 等宽 (代码/命令行): JetBrains Mono
\setmonofont{JetBrainsMono-Regular.ttf}[
    Path = font/,
    BoldFont = JetBrainsMono-Bold.ttf,
    ItalicFont = JetBrainsMono-Italic.ttf,
    BoldItalicFont = JetBrainsMono-BoldItalic.ttf]

%----- 数学字体: STIX Two Math (取自 font/ 目录, 经 unicode-math 接入) -----
% STIX Two 没有粗体数学字形, 故指定同源字体 XITS Math Bold 充当粗体
% (供向量、矩阵等 \boldsymbol 使用; 如不需要粗体数学, 可去掉 BoldFont 选项)
\RequirePackage{unicode-math}                   % OpenType 数学排版引擎
\setmathfont{STIXTwoMath-Regular.ttf}[
    Path = font/, BoldFont = XITSMath-Bold.otf]
% 证毕符号: 改用数学字体自带的白色方格, 避免使用会回退到 AMS 字体的 \square
\renewcommand{\qedsymbol}{\ensuremath{\mdlgwhtsquare}}

% 如需更换字体: 把字体文件放入 font/ 目录 (或其他自建子目录), 再修改上面相应设置即可;
% 模板只从上述路径读取字体, 不会调用 font/ 之外的任何字体.

% ===============================================================
%  二、常用宏包
% ===============================================================
% 说明: 模板所需宏包大都已由文类 boyogalazhcnbook.cls 载入, 无需在此重复:
%   数学: amsmath、amsthm、mathtools、amssymb (数学字体由第一节经 unicode-math 配置);
%   图表: graphicx、caption、tabularx (含 L/C/R/P 列格式)、booktabs、longtable、multirow、makecell、float、multicol;
%   算法: algorithm、algpseudocode; 代码: listings (样式已由文类配置, 可直接使用 lstlisting 环境).
% 下面只列出需要在本文件中载入的宏包; 用到其它宏包时在此追加即可.

\usepackage{subcaption}  % 子图排版 (caption 宏包的配套子图宏包)

%----- 绘图与流程图 (tikz) -----
% 数学书籍中的交换图、示意图、算法流程图等均可直接用 tikzpicture 环境绘制 (示例见第 2 章"算法环境"一节);
% 已载入常用库: positioning (相对定位)、arrows.meta (箭头)、shapes.geometric (菱形等流程图节点)、calc (坐标计算)
\usepackage{tikz}
\usetikzlibrary{positioning,arrows.meta,shapes.geometric,calc}

%\usepackage[notcite,notref]{showkeys}  % 调试用: 页边显示引用/标签键名

%----- 取消链接颜色和方框 (hyperref 在导言末尾载入, 设置须放入 \AtEndPreamble) -----
%\AtEndPreamble{\hypersetup{hidelinks}}

% ===============================================================
%  三、参考文献 / 引用 / 索引
% ===============================================================
%----- 参考文献格式 -----
\bibliographystyle{gbt7714-2025-numeric}     % GB/T 7714-2025 顺序编码制
%\bibliographystyle{gbt7714-2025-authoryear} % GB/T 7714-2025 著者-出版年制
%\bibliographystyle{plain}  % 其他可选: abbrv, unsrt, siam

%----- 参考文献引用格式 -----
\usepackage[numbers,sort&compress]{natbib}   % 顺序编码制引用
%\usepackage[numbers,super,square,sort&compress]{natbib}  % 上标引用
%\usepackage[authoryear,sort&compress]{natbib}            % 著者-出版年制
%\bibpunct{（}{）}{；}{a}{}{，}  % 著者-出版年制如需中文全角括号, 取消本行注释 (须配合 authoryear 引用格式)
\renewcommand{\bibfont}{\small}              % 参考文献正文字号
\setlength{\bibsep}{7pt plus 3pt minus 3pt}  % 参考文献条目间距

%----- 制作索引 -----
\usepackage{imakeidx}
\makeindex
\indexsetup{level=\chapter*}  % 索引标题的级别与章标题一致
% 提示: 编译环境若禁用了 shell 执行, makeindex 无法自动运行, 索引将为空
\ifnum\shellescape=0\relax
    \PackageWarningNoLine{boyogalazhcnbook}{编译时禁用了 shell 执行, 索引将不会自动生成;^^J请允许 shell 执行(TeX Live 默认即允许运行 makeindex), 或手动运行 makeindex main.idx}
\fi

% ===============================================================
%  四、按需微调 (取消注释后生效)
% ===============================================================
%\setlength{\itemsep}{3pt}           % 列表项之间的额外间距
%\raggedbottom                       % 允许页面底部参差, 避免出现过大的竖直空白
%\includepdf[pages=-]{pdfname.pdf}   % 插入外部 PDF 文件 (如自行制作的封面)
%\includeonly{chapters/chap03}       % 增量编译: 只编译改动过的章节 (首次需完整编译一次)
